\section{P vs NP via Quantale-Weakness and the $K_{\mathrm{poly}}$ Functor}
\label{sec:goertzel-pnp-state}

This section is a research-program lens on Goertzel's route toward separating
P and NP through observational \emph{quantale weakness} composed with a
polynomial-functor compression target ($K_{\mathrm{poly}}$). It is written to
direct creative effort. The route does \emph{not} claim P~$\neq$~NP; what is
machine-checked is a precisely localized \emph{conditional} structure --- an
obstruction packet --- plus several proved boundaries showing why the naive
shortcuts cannot work. Every conjecture below is labelled. Lean identifiers cite
the live \texttt{Computability/PNP} development.

\subsection{The creative bet}

The bet has two halves that must meet in the middle.

\emph{First, the quantale-weakness bridge.} Operational distinguishability ---
the observational calculus of which environments a machine can tell apart ---
is lifted into a quantale lattice encoding the observational closure of an
environment. The hope is that ``how weak'' a polynomial-time learner is at
distinguishing states becomes an algebraic invariant. The hard lesson already
formalized: a predictor that factors through a \emph{single global} weakness
scalar collapses to a constant
(\verb|FactorsThroughDistinctionWeakness|,
\verb|not_surjective_familyFactorsThroughDistinctionWeakness_of_nontrivial|,
\verb|GlobalWeaknessObstruction|). So global weakness data alone is provably
insufficient; the route needs richer per-state observable semantics.

\emph{Second, the $K_{\mathrm{poly}}$ compression target.} After a switching
construction the inputs factor through an exact visible surface
$u=(z,a_i,b)$ (\verb|ExactVisiblePostSwitchSurface|,
\verb|ActualSwitchedLocalInterface|). The claim that would close the route is
that the family of local Boolean rules on that surface admits a finite predictor
cover of size at most $2^s$ for polynomial $s$ ---
\verb|ExactVisibleCompressionTarget| $\equiv$ \verb|HasBitBudget| $\equiv$
\verb|FinitePredictorCover (2^s)|, an equivalence proved by
\verb|clockedKpolyFiniteLearningPayload_iff_finitePredictorCover| and
\verb|exactVisibleCompressionTarget_iff_finitePredictorCover|. If the switched
family really is so compressible, the clocked finite-learning payload exists and,
fed into the ERM recovery backbone, drives a contradiction with the assumed
fast algorithm.

The genius of the framing is its honesty about the barrier wall: instead of an
unconditional lower bound (which would crash into relativization, natural
proofs, and algebrization), the route isolates the \emph{exact minimal floor}
that the switching strategy needs --- and proves, separately, that without that
floor there is no route.

\subsection{The obstruction packet (proved, but CONDITIONAL)}

What is machine-checked is a two-way conditional, not an impossibility:
\begin{quote}
\emph{Route $\Longrightarrow$ necessary floors; and missing floor
$\Longrightarrow$ no route.}
\end{quote}
Concretely the packet proves that several tempting shortcuts \emph{cannot}
supply the floor by themselves:
\begin{itemize}\itemsep1pt
  \item locality alone (output $=$ a local rule on the visible input) permits
    the full Boolean rule class, so it cannot force a small cover ---
    \verb|fullRuleActualSwitchedLocalInterface_not_clockedKpolyFiniteLearningPayload_of_lt_surfaceCard|,
    \verb|actualSwitchedLocalInterface_not_forall_clockedKpolyFiniteLearningPayload_of_lt_surfaceCard|;
  \item plugin / hash-table lookup endpoints are still full-rule expressive ---
    \verb|pluginLookupActualSwitchedLocalInterface_not_finitePredictorCover_of_lt_surfaceCard|;
  \item a fully expressive family below the truth-table threshold has no payload
    --- \verb|not_clockedKpolyFiniteLearningPayload_of_surjective_predict_of_lt_predictorCard|;
  \item balanced-fiber source erasure neutralizes every output predicate, so it
    cannot carry compression ---
    \verb|finiteCoinRecodingFiberBalanced_iff_outputPredicateErasing|,
    \verb|finiteCoinOutputPredicateNeutral_iff_countIndependentWithin_trueSource|.
\end{itemize}
Each result peels away one false hope and sharpens what the genuine floor must
provide. Two \emph{sufficient} closure conditions are also formalized as the
floor's positive face: \verb|BitCodeData| (a bounded per-index code) and
\verb|ZABDecisionListData| (realization by a shared
$(\mathrm{zfeat}(z),a,b)$ decision-list family), with explicit budgets
$r+2k+1$ (exact ZAB) and the affine/sparse variants, closed by
\verb|exactZABCodeConstruction| and \verb|exactZABDecisionListERMConstruction|.

\subsection{The open crux: are the floors attainable?}

The route reduces to one question: \emph{is the actual switched local family
realized by a polynomial-budget code?} Equivalently, does
\verb|HasBitBudget G s| hold for the real switched family on
$u=(z,a_i,b)$, with $s$ polynomial? Two concrete roads are open:
\begin{enumerate}\itemsep1pt
  \item \textbf{ERM equality.} Show the switched family equals
    \verb|exactZABDecisionListERMFamily (zfeat, samples)| for some feature
    extractor $\mathrm{zfeat}:Z\to\mathrm{BitVec}\,r$
    (\verb|clockedKpolyFiniteLearningPayload_of_eq_exactZABDecisionListERMFamily|).
  \item \textbf{Code witness.} Supply explicit per-index codes
    $\mathrm{codes}:I\to \verb|SharedAffineDecisionListCode|(r+2k)$ realizing
    the local rules.
\end{enumerate}
This crux is \emph{exactly where the barrier wall stands}. Relativization,
natural proofs, and algebrization all forbid certain generic routes to
lower bounds. The packet's value is that it does not pretend to walk through the
wall: it stays finite and finite-count
(\verb|CountIndependentWithinDefect|,
\verb|countIndependentWithinTolerance_iff_defect_le|), works per-state via the
local visible input rather than a global uniformity argument, and uses concrete
ERM / decision-list constructions rather than diagonalization. The open question
is whether the switched family's structure --- which is dictated by the NP
problem and the switching construction, not chosen freely --- happens to land
inside the polynomial-budget cover. If it does, the route closes; if it
provably does not, that non-realizability is itself a structured lower-bound
statement.

\subsection{Conjectured workarounds for the barrier wall}

\begin{conjecture}[Speculative --- Cross-completion message rigidity]
\label{conj:pnp-rigidity}
There is a ``message rigidity'' invariant of the switched family: across the
completions of a hard NP instance, the visible surface $u=(z,a_i,b)$ must carry
strictly more than $s$ bits of mutual information with the output for any
polynomial $s$. If formalized, rigidity is precisely the statement that the
floor is \emph{not} attainable --- and, crucially, it is a property that does
\emph{not} relativize, because it is computed from the concrete combinatorics of
the switching construction, not from oracle-blind machine behavior.
\end{conjecture}

\noindent\emph{Why it might evade the wall.} The barriers bite against proof
techniques that are insensitive to the internal structure of the family
(relativizing or naturalizable). The finite-count independence machinery already
in Lean is family-specific: \verb|CountIndependentWithin| is a counting identity
on the actual fibers, not an oracle-relative statement. Rigidity would be a
quantitative refinement of that counting. \emph{Why it might fail.} If rigidity
is a ``large'' / ``constructive'' property in the natural-proofs sense, it
re-enters the wall; the conjecture is only responsible if the rigidity invariant
is shown to be neither natural (it must fail on a random family) nor
relativizing.

\begin{conjecture}[Speculative --- Per-state quantale semantics]
\label{conj:pnp-quantale}
The global-weakness collapse
(\verb|GlobalWeaknessObstruction|) is repaired by a \emph{per-state} quantale
observable: assign to each state a quantale element recording its local
observational closure, so that the predictor factors through a tensor of local
elements rather than one global scalar. The tensor's rank lower-bounds the bit
budget, converting quantale weakness directly into a
\verb|FinitePredictorCover| obstruction.
\end{conjecture}

\noindent\emph{Why it might hold.} The quantale tensor $\otimes$ is exactly the
algebra of joint evidence; a high tensor rank is the algebraic shadow of an
incompressible family. \emph{Why it might fail.} Building the per-state semantics
without secretly assuming the answer (high rank) is the whole difficulty; the
conjecture is a target, not a method.

\subsection{Where to point creative effort next}

The single highest-leverage target is
\textbf{Conjecture~\ref{conj:pnp-rigidity} (cross-completion message rigidity)},
attacked through the existing finite-count independence layer. The reason: the
obstruction packet has already shown that \emph{every} structure-blind shortcut
fails, which means the remaining live information must be a structure-\emph{aware}
invariant of the switched family --- and rigidity is the most natural candidate
that is both finite (so it is checkable against the existing
\verb|CountIndependentWithin| API) and plausibly non-relativizing (so it has a
chance against the barrier wall). The recommended push: formalize a candidate
rigidity quantity as a counting lower bound on
\verb|ExactVisiblePostSwitchSurface|, test it on small switched families where
the floor is known to be unattainable, and check whether it survives the
random-family stress test that natural proofs would demand it fail. A positive
result would be the first floor-unattainability witness; a negative result would
sharply locate which barrier the route actually hits.

\subsection*{PLN lens}

\begin{table}[h]
\centering
\small
\begin{tabular}{lcccc}
\toprule
Central node & STV $\langle s,c\rangle$ & ITV $[L,U]$ & Progress & Status \\
\midrule
Floor attainable ($K_{\mathrm{poly}}$ cover for real family) & $\langle 0.06, 0.20\rangle$ & $[0.012,\,0.812]$ & --- & open crux \\
Cross-completion message rigidity & $\langle 0.10, 0.18\rangle$ & $[0.018,\,0.838]$ & --- & open crux \\
Obstruction packet (conditional) & $\langle 0.92, 0.80\rangle$ & $[0.736,\,0.936]$ & --- & proved (conditional) \\
\midrule
\textbf{Route (through)} & $\langle 0.04, 0.23\rangle$ & $[0.0092,\,0.7792]$ & $\sim$45\% & \\
\bottomrule
\end{tabular}
\caption{PLN summary for the P-vs-NP route. The obstruction packet is
high-strength because it is a proved conditional; the through-route is very
low-strength and low-confidence because attainability of the floor sits squarely
on the relativization / natural-proofs / algebrization barrier wall, where wide
indefinite intervals are honest.}
\end{table}

\noindent Through-route: $\langle 0.04, 0.23\rangle$, ITV
$[0.04\cdot 0.23,\;0.04\cdot 0.23+(1-0.23)] = [0.0092,\,0.7792]$,
PROGRESS $\sim 45\%$. The very low strength is deliberate intellectual honesty:
a separation through this route would resolve a problem guarded by three formal
barriers, and nothing yet shows the floor is attainable or unattainable. The
wide interval encodes that we genuinely do not know which side of the wall the
switched family lies on. Progress is moderate because the conditional structure
and its boundaries are thoroughly mapped, even though the decisive crux is
untouched.
